% ============================================================
% PREAMBLE  --  shared by desk.tex, book.tex and every web PDF.
% This is the file to edit to change how everything looks.
% Compile with XeLaTeX (Overleaf: Menu > Compiler > XeLaTeX).
% ============================================================


% ============================================================
% FONTS
% ============================================================

\usepackage{fontspec}
% CMU Serif if available, otherwise FreeSerif.
\IfFontExistsTF{cmunrm.otf}{%
  \setmainfont{cmunrm.otf}[
    BoldFont=cmunbx.otf,
    ItalicFont=cmunti.otf,
    BoldItalicFont=cmunbi.otf
  ]%
}{%
  \setmainfont{FreeSerif}%
}


% ============================================================
% PAGE
% ============================================================

\usepackage{geometry}
% \pageshift moves the header and the text down together.
\newlength{\pageshift}
\setlength{\pageshift}{7.5pt}

\geometry{
  left=0.9in,
  right=1in,
  bottom=1in,
  top=\dimexpr 2cm+\pageshift\relax,
  headheight=20pt,
  headsep=15pt
}

\renewcommand{\baselinestretch}{1.145}
\setlength{\parindent}{20pt}
\setlength{\parskip}{0.55\baselineskip}

\clubpenalty=10000
\widowpenalty=10000
\displaywidowpenalty=10000


% ============================================================
% MATHEMATICS
% ============================================================

\usepackage{amssymb,amsfonts,amsmath,amsthm}
\usepackage{mathtools}
\usepackage{mathrsfs}
\usepackage[mathscr]{euscript}
\usepackage{bm}
\usepackage{bbm}
\usepackage{tensor}
\usepackage{slashed}
\usepackage{braket}

\AtBeginDocument{\let\hbar\hslash}


% ============================================================
% GENERAL PACKAGES
% ============================================================

\usepackage{graphicx}
\usepackage{microtype}
\usepackage[shortlabels]{enumitem}   % also accepts \begin{enumerate}[(a)]
\usepackage{changepage}
\usepackage{xparse}
\usepackage{etoolbox}
\usepackage{comment}
\usepackage[explicit]{titlesec}
\usepackage{titletoc}
\usepackage{multicol}
\usepackage[most]{tcolorbox}
\usepackage{caption}
\usepackage{ragged2e}
\usepackage{fancyhdr}

\usepackage[dvipsnames]{xcolor}
\definecolor{linkblue}{RGB}{28,86,140}

\usepackage[
  colorlinks=true,
  urlcolor=linkblue,
  citecolor=linkblue,
  linkcolor=linkblue,
  filecolor=linkblue,
  linktocpage=true
]{hyperref}

\usepackage{nameref}
\usepackage{cleveref}


% ============================================================
% MODE SWITCH
% book.tex says \booktrue; everything else is a standalone document.
% ============================================================

\newif\ifbook


% ============================================================
% HEADINGS
%
% \section and \subsection keep their numbered format. In small
% documents write \section*{...} and \subsection*{...}; remove the
% stars later (in a big document) and they take on numbers.
% ============================================================

% Header sizes:
%   \sectionsizefactor       section size = factor x \Large
%   \subsectionsizefraction  subsection size, as a fraction of the
%                            way from \normalsize (0) to \large (1)
\newcommand{\sectionsizefactor}{0.9}
\newcommand{\subsectionsizefraction}{0.0}

\makeatletter
\newcommand{\notes@sizeof}[2]{\begingroup#2\xdef#1{\f@size}\endgroup}
\newcommand{\notes@setsize}[1]{%
  \fontsize{#1}{\dimexpr 1.3\dimexpr #1\relax\relax}\selectfont
}
\newcommand{\sectionfont}{%
  \notes@sizeof\notes@L{\Large}%
  \notes@setsize{\dimexpr \sectionsizefactor\dimexpr \notes@L pt\relax\relax}%
}
\newcommand{\subsectionfont}{%
  \notes@sizeof\notes@n{\normalsize}%
  \notes@sizeof\notes@l{\large}%
  \notes@setsize{\dimexpr \notes@n pt
    + \subsectionsizefraction\dimexpr \notes@l pt - \notes@n pt\relax\relax}%
}
\makeatother

\newcommand{\sectionnumberbox}{\thesection\hspace{1em}}

\titleformat{\section}[hang]
  {\sectionfont\bfseries}
  {\sectionnumberbox}
  {0pt}
  {#1}

\titleformat{\subsection}[hang]
  {\normalfont\subsectionfont\bfseries}
  {\thesubsection\hspace{1em}}
  {0pt}
  {#1}
\titlespacing*{\subsection}{0pt}{0.65\baselineskip}{0.15\baselineskip}

\makeatletter
\newif\ifnotes@firstpara
\notes@firstparafalse

\let\originalsection\section
\RenewDocumentCommand{\section}{s o m}{%
  \closestatement
  \IfBooleanTF{#1}{%
    \originalsection*{#3}%
  }{%
    \IfNoValueTF{#2}{\originalsection{#3}}{\originalsection[#2]{#3}}%
  }%
  \global\notes@firstparatrue
}

\let\originalsubsection\subsection
\RenewDocumentCommand{\subsection}{s o m}{%
  \closestatement
  \IfBooleanTF{#1}{%
    \originalsubsection*{#3}%
  }{%
    \IfNoValueTF{#2}{\originalsubsection{#3}}{\originalsubsection[#2]{#3}}%
  }%
  \global\notes@firstparatrue
}
\makeatother


% ============================================================
% DOCUMENTS
%
%   \smalldoc{Title}{Tag}  first line of every document. Prints the title
%                          like \section*{Title} (in the book: a numbered
%                          \section). The tag (Cov, Mod, SUSY) names the
%                          document for the web and for \docref.
%   \uses{Tag1, Tag2}      lines on the web (prints nothing)
%   \near{Tag}             where a new node first appears on the web
%                          (prints nothing)
%   \docref{Tag}           link to another document in the text
% ============================================================

\newcommand{\doctag}{}

\makeatletter
\newcommand{\smalldoc}[2]{%
  \closestatement
  \renewcommand{\doctag}{#2}%
  \ifbook\section{#1}\else\section*{#1}\fi
  \phantomsection\label{doc:#2}%
  \setcounter{equation}{0}\setcounter{statement}{0}%
  \global\notes@firstparatrue
}
\makeatother

\newcommand{\docref}[1]{%
  \ifbook\hyperref[doc:#1]{\textsf{#1}}%
  \else\csname hyper@linkurl\endcsname{\textsf{#1}}{#1.pdf}\fi}

% read by the web build only; they print nothing
\newcommand{\uses}[1]{}
\newcommand{\near}[1]{}
\newcommand{\assumes}[1]{}
\newcommand{\status}[1]{}

% Equations: (1), (2), ... in each document; (3.1), (3.2), ... in the book.
\renewcommand{\theequation}{\ifbook\thesection.\fi\arabic{equation}}
\renewcommand{\theHequation}{\doctag.\arabic{equation}}


% ============================================================
% THE BOOK  (book.tex)
%
%   \documentheading{TITLE}{Author \\ Date}
%   \tableofcontents          two columns, sections only
%   \input{docs/...}          each document becomes a numbered section
%   \stubsection{Title}       planned section: "--" for its page in the
%                             contents, nothing in the text
% ============================================================

\newcommand{\documentheading}[2]{%
  \begin{center}
    {\large\bfseries\MakeUppercase{#1}}%
    \par\vspace{1.55\baselineskip}%
    #2
  \end{center}%
  \vspace{1.55\baselineskip}%
}

\setcounter{tocdepth}{1}

\newlength{\tocnumwidth}
\settowidth{\tocnumwidth}{105\quad}
\newcommand{\tocpagebox}[1]{\makebox[1.7em][r]{#1}}

\titlecontents{section}
  [0pt]
  {\addvspace{0.15em}\normalfont}
  {\makebox[\tocnumwidth][l]{\thecontentslabel}}
  {}
  {\normalfont\ \titlerule*[0.6pc]{.}\ \tocpagebox{\thecontentspage}}

\titlecontents{stubsection}
  [0pt]
  {\addvspace{0.15em}\normalfont}
  {\makebox[\tocnumwidth][l]{\thecontentslabel}}
  {}
  {\normalfont\ \titlerule*[0.6pc]{.}\ \tocpagebox{--}}

\contentsmargin{0pt}

\makeatletter
\def\toclevel@stubsection{1}
\renewcommand{\tableofcontents}{%
  \begingroup
  \small
  \linespread{0.88}\selectfont
  \setlength{\columnsep}{2.5em}%
  \parindent=0pt
  \parskip=1pt
  \begin{multicols}{2}%
  \raggedcolumns
  \@starttoc{toc}%
  \end{multicols}%
  \endgroup
}
\newcommand{\stubsection}[1]{%
  \refstepcounter{section}%
  \addcontentsline{toc}{stubsection}{\protect\numberline{\thesection}#1}%
}
\makeatother


% ============================================================
% TEXT UNITS
%
%   \paragraph{Groups} Text...    ->  Groups.  Text...   bold run-in title
%   \paragraph{} Text...          ->  plain paragraph break
%   \paragraph{Example} Text...   ->  Example 2.  Text...  (a statement)
%
% Titles registered with \statementtitle give statements; every
% other title is a bold run-in title.  Statements can also be
% written directly, with an optional name:
%
%   \definition  \axiom  \proposition[path lifting]  \theorem
%   \lemma  \corollary  \fact  \example  \remark  \exercise
%   and their plurals (\examples, \corollaries, ...), plus \proof
%
% \cref{label} gives "§3.2" for a numbered section and "Proposition 3" for a statement.
% ============================================================

\newlength{\factgap}
\setlength{\factgap}{0.7em}

% One counter for all statements in a document:
% Definition 1, Example 2, ... in each document; Definition 3.1, ... in the book.
\newcounter{statement}
\renewcommand{\thestatement}{\ifbook\thesection.\fi\arabic{statement}}
\renewcommand{\theHstatement}{\doctag.\arabic{statement}}

% Label style for statements.  Choose one:
%   \statementstyle{italic}   Proposition 2 (name).  Text...
%   \statementstyle{margin}   label small and grey in the left margin
\newcommand{\statementstyle}[1]{\def\notesstatementstyle{#1}}
\statementstyle{italic}

% A positive \statementindent sets statements narrower on both sides.
\newlength{\statementindent}
\setlength{\statementindent}{0pt}

\newcommand{\closestatement}{%
  \par
  \global\leftskip=0pt
  \global\rightskip=0pt
}
\newcommand{\openstatement}{%
  \closestatement
  \addvspace{0.55\baselineskip}%
  \global\leftskip=\statementindent
  \global\rightskip=\statementindent
}

\newcommand{\marginlabelfont}{\footnotesize\color{black!55}}
\newlength{\marginlabelsep}
\setlength{\marginlabelsep}{9pt}

% \statement[name]{Label}{Abbreviation}{cref type}
\NewDocumentCommand{\statement}{o m m m}{%
  \openstatement
  \crefalias{statement}{#4}%
  \refstepcounter{statement}%
  \noindent
  \ifdefstring{\notesstatementstyle}{margin}{%
    \llap{\marginlabelfont #3~\thestatement\hspace{\marginlabelsep}}%
    \IfValueT{#1}{{\bfseries\MakeUppercase#1.}\hspace{\factgap}}%
  }{%
    {\bfseries #2~\thestatement}%
    \IfValueT{#1}{ (#1)}%
    {\bfseries .}%
    \hspace{\factgap}%
  }%
  \ignorespaces
}

% Singular and plural forms; both take an optional name.
\NewDocumentCommand{\definition}{o}{\statement[#1]{Definition}{Def.}{definition}}
\NewDocumentCommand{\definitions}{o}{\statement[#1]{Definitions}{Def.}{definition}}
\NewDocumentCommand{\axiom}{o}{\statement[#1]{Axiom}{Ax.}{axiom}}
\NewDocumentCommand{\axioms}{o}{\statement[#1]{Axioms}{Ax.}{axiom}}
\NewDocumentCommand{\proposition}{o}{\statement[#1]{Proposition}{Prop.}{proposition}}
\NewDocumentCommand{\propositions}{o}{\statement[#1]{Propositions}{Prop.}{proposition}}
\NewDocumentCommand{\theorem}{o}{\statement[#1]{Theorem}{Thm.}{theorem}}
\NewDocumentCommand{\theorems}{o}{\statement[#1]{Theorems}{Thm.}{theorem}}
\NewDocumentCommand{\lemma}{o}{\statement[#1]{Lemma}{Lem.}{lemma}}
\NewDocumentCommand{\lemmas}{o}{\statement[#1]{Lemmas}{Lem.}{lemma}}
\NewDocumentCommand{\corollary}{o}{\statement[#1]{Corollary}{Cor.}{corollary}}
\NewDocumentCommand{\corollaries}{o}{\statement[#1]{Corollaries}{Cor.}{corollary}}
\NewDocumentCommand{\fact}{o}{\statement[#1]{Fact}{Fact}{fact}}
\NewDocumentCommand{\facts}{o}{\statement[#1]{Facts}{Fact}{fact}}
\NewDocumentCommand{\example}{o}{\statement[#1]{Example}{Ex.}{example}}
\NewDocumentCommand{\examples}{o}{\statement[#1]{Examples}{Ex.}{example}}
\NewDocumentCommand{\remark}{o}{\statement[#1]{Remark}{Rem.}{remark}}
\NewDocumentCommand{\remarks}{o}{\statement[#1]{Remarks}{Rem.}{remark}}
\NewDocumentCommand{\exercise}{o}{\statement[#1]{Exercise}{Exer.}{exercise}}
\NewDocumentCommand{\exercises}{o}{\statement[#1]{Exercises}{Exer.}{exercise}}

% Proofs are not numbered.
\RenewDocumentCommand{\proof}{}{%
  \openstatement
  \noindent{\itshape Proof.}\hspace{\factgap}\ignorespaces
}

\makeatletter
% If ordinary prose comes first after a heading, clear the flag when that paragraph ends.
\AddToHook{para/end}{\global\notes@firstparafalse}
% \statementtitle{Title}{\command}: a \paragraph{Title} produces
% that statement instead of a bold run-in title.
\newcommand{\statementtitle}[2]{%
  \expandafter\def\csname notes@title@\detokenize{#1}\endcsname{#2}%
}
% \paragraph{Title} Text: the space after the closing brace is ignored,
% so writing "\paragraph{Title} Text" and "\paragraph{Title}Text" look the same.
\renewcommand{\paragraph}[1]{%
  \closestatement
  \let\notes@next\ignorespaces
  \if\relax\detokenize{#1}\relax
    % The first paragraph after a heading is flush left; others indent.
    \ifnotes@firstpara
      \noindent
    \else
      \noindent\hspace*{\parindent}%
    \fi
  \else
    \ifcsname notes@title@\detokenize{#1}\endcsname
      \expandafter\let\expandafter\notes@next\csname notes@title@\detokenize{#1}\endcsname
    \else
      \noindent{\bfseries #1.}\hspace{\factgap}%
    \fi
  \fi
  \global\notes@firstparafalse
  \notes@next
}
\makeatother

% Titles of \paragraph{...} that produce a statement.
\statementtitle{Definition}{\definition}
\statementtitle{Definitions}{\definitions}
\statementtitle{Axiom}{\axiom}
\statementtitle{Axioms}{\axioms}
\statementtitle{Proposition}{\proposition}
\statementtitle{Propositions}{\propositions}
\statementtitle{Theorem}{\theorem}
\statementtitle{Theorems}{\theorems}
\statementtitle{Lemma}{\lemma}
\statementtitle{Lemmas}{\lemmas}
\statementtitle{Corollary}{\corollary}
\statementtitle{Corollaries}{\corollaries}
\statementtitle{Fact}{\fact}
\statementtitle{Facts}{\facts}
\statementtitle{Example}{\example}
\statementtitle{Examples}{\examples}
\statementtitle{Remark}{\remark}
\statementtitle{Remarks}{\remarks}
\statementtitle{Exercise}{\exercise}
\statementtitle{Exercises}{\exercises}
\statementtitle{Proof}{\proof}

\newcommand{\stars}{%
  \closestatement
  \addvspace{1.15\baselineskip}%
  \begin{center}%
    \textasteriskcentered\hspace{0.45em}%
    \textasteriskcentered\hspace{0.45em}%
    \textasteriskcentered\hspace{0.45em}%
    \textasteriskcentered\hspace{0.45em}%
    \textasteriskcentered
  \end{center}%
  \addvspace{1.15\baselineskip}%
}

\crefformat{section}{#2$\S$#1#3}
\crefformat{subsection}{#2$\S$#1#3}
\crefname{definition}{Definition}{Definitions}
\crefname{axiom}{Axiom}{Axioms}
\crefname{proposition}{Proposition}{Propositions}
\crefname{theorem}{Theorem}{Theorems}
\crefname{lemma}{Lemma}{Lemmas}
\crefname{corollary}{Corollary}{Corollaries}
\crefname{fact}{Fact}{Facts}
\crefname{example}{Example}{Examples}
\crefname{remark}{Remark}{Remarks}
\crefname{exercise}{Exercise}{Exercises}


% ============================================================
% PAGE NUMBERS
% Bottom-centered on every page; no running headers.
% ============================================================

\pagestyle{fancy}
\fancyhf{}
\fancyfoot[C]{$-$ \thepage\  $-$}
\renewcommand{\headrulewidth}{0pt}
\renewcommand{\footrulewidth}{0pt}

\fancypagestyle{plain}{%
  \fancyhf{}%
  \fancyfoot[C]{$-$ \thepage\  $-$}%
  \renewcommand{\headrulewidth}{0pt}%
  \renewcommand{\footrulewidth}{0pt}%
}


% ============================================================
% BOXES
%   \begin{summarybox} ... \end{summarybox}   (also: mainidea)
% ============================================================

\newenvironment{summarybox}{%
  \closestatement
  \begin{tcolorbox}[
    enhanced, breakable,
    colback=white, colframe=black,
    boxrule=0.5pt, arc=0pt, outer arc=0pt,
    left=9pt, right=9pt, top=7pt, bottom=7pt,
    before skip=\baselineskip,
    after skip=\baselineskip,
    fontupper=\normalfont
  ]%
}{%
  \end{tcolorbox}%
}
\newenvironment{mainidea}{\begin{summarybox}}{\end{summarybox}}


% ============================================================
% FIGURES AND TABLES
% ============================================================

\DeclareCaptionFormat{custom}{\small{#1#2}\justifying #3}
\captionsetup[figure]{format=custom, name=FIG., labelsep=period}
\captionsetup[table]{format=custom, name=TAB., labelsep=period}


% ============================================================
% YOUR OWN MACROS BELOW
% ============================================================
